\chapter{位置很多时，为什么还要补一小段}\label{sec:wide}
\section{看上去接近4p，实际还差一个条件}
整数下降证明的是
\[n\ge4p-2+d_p(n).
\]
它不是无条件的 $n\ge4p$。例如 $p=761,n=3045=4p+1$，而 $n=3\cdot5\cdot7\cdot29$ 有四个危险素因子。右边要求 $n\ge3046$，差一个位置。

所以即使临界带 $2p<n<4p$ 已全部完成，也不能忘了这个窄缝。我们需要证明宽维数中被旧不等式漏掉的部分也成立。

\section{一个失败的不等式，反而告诉我们有更多资源}
旧结论已经覆盖所有 $n\ge4p+\lceil\log_2p\rceil$。因此漏掉的只能写成 $n=4p+s$，其中 $s$ 很小。原不等式不成立，又意味着
\[\omega(n)\ge s+3.
\]
也就是说，位置数含很多\emph{不同}素因子。不同素数的乘积长得很快，于是 $n$、进而 $p$ 必须相当大。这给出 $p\ge16s+21$ 等强估计。

这一步改变了我们看剩余参数的方式。不是处理任意 $p,s$，而是处理“$s$ 小，$p$ 被素因子条件逼得很大”的参数。它们有足够的实际副本做一组五步反射；同时欧拉函数足够大，使安全数字有足够密度。

\section{这一讲的七步路线}
先用同样的见证集中和三行折叠，得到加权核心
\[(a^{2p},(-2a+rz)^p,(-pz)^r),\qquad r=p+s.
\]
这次第一块是 $2p$ 份，不能把它当作两个可随意交换的 $p$ 块去套前一讲的对称公式。我们重新写两套真实返回：一套五步反射和一套保留旧载体的三步反射。比较它们，取得双向剪切和明确同余层。

接下来，来源计数允许一整段安全整数 $q$。其对应的行列式成为环中单位。只要这段及其负数包含超过一半的模 $n$ 单位，每个单位就是两者之比：集合 $S$ 与 $uS$ 都超过半数，必相交，交点给 $u=a/b$。

最后一个返回提供局部平移。模五处的步长可能消失，所以先用另一返回处理五进分量，再用指数CRT处理其余素数而不动刚做好的部分。完整列匹配后，精确送到 $(\epsilon^{2p},(-2\epsilon)^p,0^{p+s})$，在其中的 $4p$ 元零和子系统收尾。

下面保留七步的全部计数。第一次阅读可以先对照这条路线确认每组公式承担哪项任务，再第二遍检查上下界。
\InsertProof{prop:wide}

\section{为什么这个收尾没有绕圈}
最终调用的是 $4p$ 元判据。它早在整数下降中独立完成，因为 $d_p(4p)=1$，故 $4p\ge4p-2+1$。我们不是用“所有大于等于 $4p$ 已解决”来证明“所有大于等于 $4p$ 已解决”，而是用其中那个先已证明的固定维数 $4p$ 消去所有非零位置。

这种先后关系值得在整份证明里一直保持：先证明一个含零或等块的小工具，后面的矩阵构造只需把状态送进它能处理的范围。
